\documentclass[tikz,border=12pt]{standalone}
\usepackage[fino]{diagrama}
% Un árbol de tipos: x del centro, y de la raíz, raíz, tres hijos
\newcommand\arbol[6]{%
  \node[font=\bfseries\px{33}] at (#1-0.08,#2+1.53) {Type Level};
  \node at (#1,#2+0.88) {T1};
  \node at (#1,#2) {#3};
  \draw[lleno] (#1,#2-0.37) -- (#1,#2-1.67);
  \draw[lleno, rounded corners=6pt] (#1,#2-1.04) -- (#1-1.8,#2-1.04) -- (#1-1.8,#2-1.67);
  \draw[lleno, rounded corners=6pt] (#1,#2-1.04) -- (#1+1.8,#2-1.04) -- (#1+1.8,#2-1.67);
  \node at (#1-1.8,#2-2.07) {#4};
  \node at (#1,#2-2.07) {#5};
  \node at (#1+1.8,#2-2.07) {#6};
  \node at (#1,#2-2.89) {T0};
  \node[font=\bfseries\px{33}] at (#1-0.08,#2-3.72) {Type Level};
}
\begin{document}
\begin{tikzpicture}[color=gris, every node/.append style={font=\px{33}, text=tinta}, line width=1.1pt,
    every picture/.append style={}, >={Latex[length=3.6mm,width=2.6mm]}]
  \tikzset{lleno/.style={-{Latex[length=3.6mm,width=2.6mm]}}}
  \node[titulo=62] at (8.6,-0.78) {Paradoxes of Nonsense in Signification};

  \node[font=\bfseries\itshape\px{40}] at (8.85,-3.1) {Regular Law};
  \arbol{5.4}{-5.63}{Animal}{Bird}{Man}{Tick}
  \arbol{11.8}{-5.63}{Plant}{Herb}{Fern}{Fruit}
  \draw[<->] (6.25,-5.68) -- (10.98,-5.68);

  \node[font=\bfseries\itshape\px{40}] at (8.85,-11.95) {Nonsense};
  \arbol{5.73}{-14.55}{Animal}{Animal}{Man}{Tick}
  \arbol{12.13}{-14.55}{Plant}{Herb}{Fern}{Animal}
  \draw[<->] (6.6,-14.58) -- (11.3,-14.58);
  \draw[puntitos, line width=1.2pt, -{Triangle[open,length=3.6mm,width=3mm]}]
    (3.27,-16.62) .. controls (2.55,-16.6) and (2.7,-14.55) .. (4.88,-14.55);
  \draw[puntitos, line width=1.2pt, -{Triangle[open,length=3.6mm,width=3mm]}]
    (14.92,-16.62) .. controls (15.65,-16.6) and (15.5,-14.55) .. (12.95,-14.55);
  \node[font=\itshape\px{33}] at (1.3,-15.65) {Paradox of the\\Abnormal Set};
  \node[font=\itshape\px{33}] at (16.7,-15.65) {Paradox of the\\Rebel Element};
\end{tikzpicture}
\end{document}
